\section{Temperature response and the passage through the small-activity region}
\label{sec:trajectory}

The activity estimates now determine the scale at which the height
variance first departs from its critical value.  The essential point is
to determine how the physical temperature crosses the critical
trajectory.  We first prove a nonzero linear response in its growing
direction, then follow a real temperature perturbation until it reaches
a fixed small exit level.  All coordinates and Gaussian reference
parameters in this section are those of Proposition~\ref{prop:map};
only the input activities depend on the physical inverse temperature
\(b\).  The comparison of changing entries follows the physical-readout
method of \cite[Lemma~5.7 and the proof of Proposition~5.1]{F};
here we apply it also to temperature derivatives and secants.

Write \(\mathcal H^{(k)}\) for the critical families in
Proposition~\ref{prop:critical}, with observation side
\(s_k\asymp2^k\).  The index \(j\) is always the absolute block index.
In particular the corresponding physical torus has side
\(n_j=D_0L^j\), even when the entry side changes.  We use
\[
 W_j=(x_j,y_j),\qquad
 x_j=b_*u_j,\qquad y_j=\sqrt{b_*H}\,Z_j,
 \qquad b_*=2\log L.
\]
The symbol \(y_j\) here denotes the harmonic coordinate, rather than
an observation field.  A superscript \(\mathrm{tor}\) on a remainder
always refers to the separately initialized torus at its specified
stop.

We fix the order of the smallness choices.  Enlarge the remainder
contraction constant in Proposition~\ref{prop:map} slightly to a
number \(\theta_1<1\), and choose \(r\in(\theta_1,1)\).  Choose a
constant \(C_*\) so large that
\begin{equation}\label{eq:traj-parameter-choice}
 C_*\log L>1000,\qquad
 \frac2{C_*}<\frac{-\log r}{100}.
\end{equation}
Next choose \(\gamma>0\) with
\[
 \gamma C_*<\frac1{100},\qquad
 \gamma<\frac{-\log r}{100}.
\]
Choose a sufficiently small activity tube, so that the derivative of
the plane map in balanced coordinates has norm at most
\(e^\gamma\) at every point of that tube.  This is possible because
its linearization is the identity on the two scalar coordinates and
a contraction on the remainder.  The nonlinear remainder terms are
absorbed while keeping their contraction at most \(\theta_1\).
Finally choose the entry parameter and its lower scale threshold as
in Propositions~\ref{prop:entry} and~\ref{prop:critical}, for this
tube.  Every choice in this paragraph is fixed independently of
\(b\).  Changing the equivalent coordinate norms only changes fixed
constants in the estimates below.

\subsection{The derivative of the critical trajectory}

\begin{proposition}[Nonzero temperature response]\label{prop:response}
For the fixed choices above there is a number
\(\mathcal P\in(0,\infty)\) such that, with
\[
 V_j^{(k)}=
 \left.\frac{\partial W_j^{(k)}(b)}{\partial b}\right|_{b=\bc},
 \qquad q_u=(1,-\tfrac12),
\]
we have
\begin{equation}\label{eq:traj-response-limit}
 \sup_{\frac12 2^k\le j\le\frac52 2^k}
 \left|\frac{V_j^{(k)}}j-\mathcal Pq_u\right|
 \longrightarrow0.
\end{equation}
The plane transverse remainder and every separately initialized
torus remainder, at an index in this range behind a valid stop in
the critical history, have temperature derivative bounded by
\(C|V_j^{(k)}|/j+e^{-cj}\).
\end{proposition}

\begin{proof}
We omit the history superscript until we compare two entries.  The
entry analyticity in Proposition~\ref{prop:entry} gives an initial
derivative bound \(Cs_k^2\).  Propagation in the chosen tube bounds
all plane derivatives by \(Cs_k^2e^{\gamma j}\); contraction of a
torus remainder driven by the same scalars gives the weaker common
bound \(Cs_k^2e^{2\gamma j}\).  These estimates are used only to
remove the initial remainder contribution.

On a slightly larger fixed fractional interval than that in the
statement, Proposition~\ref{prop:critical} gives
\[
 W_j=\frac{(1,1)}j+O\!\left(\frac{\log j}{j^2}\right),\qquad
 \|Q_j\|+\|\mathcal R_{j|h}^{\mathrm{tor}}\|=O(j^{-2}).
\]
Here \(h\ge j\) is any valid stop within the history's horizon.
Differentiating the map of Proposition~\ref{prop:map} therefore gives
\begin{equation}\label{eq:traj-linear-response}
 V_{j+1}=A_jV_j+B_j\partial_b Q_j,
 \qquad
 A_j=I+\frac{\mathsf M}{j}
            +O\!\left(\frac{\log j}{j^2}\right),
 \qquad \|B_j\|\le\frac Cj,
 \qquad
 \mathsf M=\begin{pmatrix}0&-2\\-1&-1\end{pmatrix}.
\end{equation}
For a fixed torus stop \(h\), put
\[
 r_{j|h}=\|\partial_bQ_j\|
                 +\|\partial_b\mathcal R_{j|h}^{\mathrm{tor}}\|.
\]
The analytic remainder estimates, the absence of linear forcing
from the scalar inputs into the unbalanced torus remainder, and the
coordinate changes in Proposition~\ref{prop:map} imply
\begin{equation}\label{eq:traj-derivative-contraction}
 r_{j+1|h}\le\theta_1r_{j|h}+\frac Cj|V_j|.
\end{equation}
Terms with a small coefficient multiplying either remainder
derivative have been included in \(\theta_1\).  The constants are
independent of \(h\).

We give the backward estimate that makes this contraction useful.
Fix \(j\) in the asserted interval and take
\(l=\lfloor j/2\rfloor\).  For large \(k\), the interval
\([l,j]\) lies inside the slightly enlarged fractional range on
which the preceding bounds hold.  Define
\[
 S_v=\max_{l\le i\le j}r^{j-i}|V_i|,
 \qquad
 S_q=\max_{l\le i\le j}r^{j-i}r_{i|h}.
\]
Iterating \eqref{eq:traj-derivative-contraction}, and using
\(i\asymp j\) on this interval, gives
\begin{equation}\label{eq:traj-weighted-forward}
 S_q\le\frac CjS_v+Cr^{j-l}r_{l|h}.
\end{equation}
For large \(j\), the matrices \(A_i\) are invertible and
\(\|A_i^{-1}\|\le1+C/j\) in a fixed equivalent scalar norm.
Solving \eqref{eq:traj-linear-response} backward from \(j\),
multiplying at index \(i\) by \(r^{j-i}\), and summing the
geometric factors \(r^{v-i}\) over intermediate indices
\(i\le v<j\), gives
\begin{equation}\label{eq:traj-weighted-backward}
 S_v\le C|V_j|+\frac CjS_q.
\end{equation}
Absorbing \(CS_q/j^2\) between these inequalities yields
\[
 r_{j|h}\le\frac Cj|V_j|+Cr^{j-l}r_{l|h}.
\]
The crude initial derivative estimate and the choice of \(\gamma\)
make the last term exponentially small on these fractional
intervals.  Consequently
\begin{equation}\label{eq:traj-slaving}
 r_{j|h}\le\frac Cj|V_j|+e^{-cj}.
\end{equation}
The same proof without the torus term gives the plane assertion.

We next use the physical meaning of the two scalar coordinates.
The height part of the first-cosine observation has variance bounded
below by a positive constant on the critical tori and has bounded
reference cost.  Lemma~\ref{lem:response} thus gives a fixed positive
lower bound for its temperature derivative.  Lemma~\ref{lem:cutoff-derivative}
allows us to use the cutoff readout with an exponentially small
error.  Differentiating Proposition~\ref{prop:readouts}, using
\eqref{eq:traj-slaving} and the coordinate changes, bounds this
response by \(C|V_j|+e^{-cj}\).  In particular
\begin{equation}\label{eq:traj-response-nonvanishing}
                         |V_j|\ge c>0.
\end{equation}
This lower bound is the step that will exclude a purely decaying
solution of the linearized recurrence.

To compare two entries, define the normalized diagnostic vector
\begin{equation}\label{eq:traj-normalized-diagnostic}
 \mathscr O_{j,s}(b)=
 \begin{pmatrix}
 \bigl(\widehat\Var_b X_{n_j,s}
                  -\|f_{n_j,s}\|_2^2\bigr)/\beta_{n_j,s}^2\\[2pt]
 e^{(s/n_j)^2/2}\widehat\E_b e^{i\bar\eta}
 \end{pmatrix}.
\end{equation}
The subtracted term is the independent observation-noise variance
in the uncut law.  Lemma~\ref{lem:grid-comparison} compares this
vector with the same physical pair
\((\Var_b^{\mathrm{tor}}H_{n_j},
\E_b^{\mathrm{tor}}e^{i\bar h})\) for either entry.
At zero activity, its linear scalar coefficient in the
\(W_j,Q_j\) coordinates is
\[
 \operatorname{diag}\left(
 \frac{\kappa_{n_j,s}}{b_*\beta_{n_j,s}^2},
 \frac{e^{(s/n_j)^2/2}p_{j,s}\ell_j^{-1}}{\sqrt{b_*H}}\right).
\]
Here \(\ell_j=1+O(e^{-cj})\) is the balancing factor in the
proof of Proposition~\ref{prop:map}.  The linear remainder rows
are controlled by \eqref{eq:traj-slaving}.
Its constant background is independent of \(b\).  Differentiating
and using \eqref{eq:traj-slaving} gives
\begin{equation}\label{eq:traj-response-readout}
 D_jV_j+O(|V_j|/j)+O(e^{-cj}),
 \qquad
 D_j=\operatorname{diag}\left(
       \frac{a\mathfrak g}{b_*},
       \frac{p_j}{\sqrt{b_*H}}\right).
\end{equation}
The numbers \(p_j\) and their inverses are uniformly bounded and
positive.  At a common absolute stop this diagonal matrix may be
chosen identically for both entries: the difference of the actual
multipliers is exponentially small and is included in the error.
The variance comparison error is bounded by
\(Cn_j^2(s/n_j)^4\); the phase comparison is exact after removal
of the mean noise, apart from the cutoff error.  All these errors
are exponentially small when the entries being compared have
\(s\asymp j\).  Since both readouts describe the same finite
physical torus, \eqref{eq:traj-response-readout} and
\eqref{eq:traj-response-nonvanishing} give
\begin{equation}\label{eq:traj-response-overlap}
 |V_j^{(k)}-V_j^{(k')}|
 \le \frac Cj\bigl(|V_j^{(k)}|+|V_j^{(k')}|\bigr)
\end{equation}
whenever the index lies in the fixed interior ranges for both
histories.  In particular their norms are comparable.  Only two
physical scalar readouts were compared; no subtraction of
remainders belonging to different activity norms is involved.

Define a single sequence by taking \(V_j=V_j^{(k)}\) for
\(2^k\le j<2^{k+1}\).  At an ordinary step,
\eqref{eq:traj-linear-response} and \eqref{eq:traj-slaving} show
that this sequence obeys
\[
 V_{j+1}=\left(I+\frac{\mathsf M}{j}+E_j\right)V_j,
 \qquad \|E_j\|\le C\frac{\log j}{j^2}.
\]
Indeed an error vector of norm at most \(\eta|V_j|\) can be
written as a matrix of norm at most \(\eta\) applied to \(V_j\).
The nonvanishing bound permits this representation also for the
exponentially small errors.  At a change of entry there is, by
\eqref{eq:traj-response-overlap}, an additional matrix error of
norm \(C/j\).

The matrix \(\mathsf M\) has eigenvectors
\(q_u=(1,-1/2)\) and \(q_s=(1,1)\), with eigenvalues \(1\) and
\(-2\), respectively.  Let \((U_k,S_k)\) be the coordinates of
\(V_{2^k}\) in this basis.  Telescoping products over a dyadic
interval gives
\begin{equation}\label{eq:traj-epoch}
 \begin{pmatrix}U_{k+1}\\S_{k+1}\end{pmatrix}
 =\left\{
   \begin{pmatrix}2&0\\0&1/4\end{pmatrix}+\mathsf E_k
   \right\}
   \begin{pmatrix}U_k\\S_k\end{pmatrix},
 \qquad \|\mathsf E_k\|\le Ck2^{-k}.
\end{equation}
For completeness, the unperturbed growing product is exactly
\(\prod_{j=n}^{2n-1}(1+1/j)=2\), and the decaying product is
\(\prod_{j=n}^{2n-1}(1-2/j)=1/4+O(n^{-1})\).
All partial products within one dyadic interval and their inverses
are bounded.  Summing \(C\log j/j^2\) over the interval, together
with its one \(C/j\) change of entry, proves the stated bound on
\(\mathsf E_k\).

For all sufficiently large \(k\), the cone
\(\{|S_k|\le|U_k|\}\) is invariant under
\eqref{eq:traj-epoch}.  If a sequence never enters this cone,
then \(|U_k|<|S_k|\) at every sufficiently large index, and
\eqref{eq:traj-epoch} gives
\(|S_{k+1}|\le\tfrac12|S_k|\).  Both coordinates would tend to
zero, contrary to \eqref{eq:traj-response-nonvanishing}.  Thus the
sequence enters the cone.  Thereafter
\[
 U_{k+1}=2U_k(1+\eta_k),\qquad
 |\eta_k|\le Ck2^{-k}.
\]
The factors \(1+\eta_k\) are positive on the tail and their
logarithms are absolutely summable.  It follows that
\(U_k/2^k\) converges to a finite nonzero number \(\mathcal P\).
The second row of \eqref{eq:traj-epoch}, divided by its first,
gives
\[
 \frac{|S_{k+1}|}{|U_{k+1}|}
 \le \rho\frac{|S_k|}{|U_k|}+Ck2^{-k}
 \quad\text{for some }\rho<1,
\]
so this ratio tends to zero.  The same product estimates on partial
dyadic intervals show that \(V_j/j\to\mathcal Pq_u\) at every
intermediate index.  Applying \eqref{eq:traj-response-overlap}
then gives the uniform assertion for either history on its interior
range.  Finally the first component of
\eqref{eq:traj-response-readout} has positive physical response.
Its leading term would be negative and of order \(j\) if
\(\mathcal P<0\).  Hence \(\mathcal P>0\).
\end{proof}

\subsection{Matching a late entry to the critical trajectory}

We will use one observation side much larger than the absolute index
at which the perturbation is first measured.  The next lemma supplies
the critical estimates for that history.  Its proof compares finite
physical diagnostics, so no rate of convergence of a critical height
scaling limit is needed.

\begin{lemma}[Critical estimates after a logarithmic initial interval]
\label{lem:traj-late-critical}
Let \(s\to\infty\) through admissible observation sides, put
\(J=\lfloor C_*\log s\rfloor\), and let \(K(s)\ge J\) satisfy
\(K(s)=o(s^2)\).  For the history with observation side \(s\), at
\(b=\bc\), uniformly for \(J/2\le i\le K(s)\),
\begin{equation}\label{eq:traj-late-critical}
 W_i=\frac{(1,1)}i+O\!\left(\frac{\log i}{i^2}\right),
 \qquad
 \|Q_i\|+\|\mathcal R_i\|=O(i^{-2}).
\end{equation}
Every separately initialized torus with a valid stop
\(h\in[i,K(s)]\) also has
\(\|\mathcal R_{i|h}^{\mathrm{tor}}\|=O(i^{-2})\).
\end{lemma}

\begin{proof}
The extended critical confinement in Proposition~\ref{prop:critical}
keeps all these states in the chosen small tube.  Put
\(i_0=\lceil J/100\rceil\).  At every \(i\ge i_0\), compare the
history with a canonical history \(\mathcal H^{(k)}\) for which
\(2^k\le i<2^{k+1}\).  Both are exact descriptions of the same
physical torus of side \(n_i\).  The first choice in
\eqref{eq:traj-parameter-choice} ensures
\(s/n_i\le e^{-ci}\) throughout this range.  The observation
sampling errors and multiplier differences are consequently
exponentially small.  The cutoff error is uniformly at most
\[
 C\exp\{CK(s)-cs^2\},
\]
and is also exponentially small in \(i\), since \(K(s)=o(s^2)\).
The canonical history has the bounds of
Proposition~\ref{prop:critical}.

Write \(d_i=|t_i|+|z_i|\), and let \(R_i\) be the maximum of
\(\|\mathcal R_i\|\) and the remainder norms at stage \(i\)
of all the separate tori with stops in \([i,K(s)]\).  The two
nonzero scalar multipliers in Proposition~\ref{prop:readouts} give
\begin{equation}\label{eq:traj-late-readout}
 d_i\le\frac Ci+CR_i+Cd_i^2.
\end{equation}
In writing this inequality we absorbed remainder squares and mixed
terms using the uniform small tube.  Remainder contraction, with
copied scalar forcing also on each torus, gives
\begin{equation}\label{eq:traj-late-contraction}
 R_{i+1}\le\theta_1R_i+Cd_i^2.
\end{equation}
The family of tori in the maximum only decreases when \(i\)
increases, so the same inequality holds for the maximum itself.

Absorb \(Cd_i^2\) into the left side of
\eqref{eq:traj-late-readout}.  Then
\(d_i\le C/i+CR_i\), and, since \(R_i\) is uniformly small,
\eqref{eq:traj-late-contraction} becomes
\[
 R_{i+1}\le\theta_2R_i+\frac C{i^2}
 \quad\text{for a fixed }\theta_2<1.
\]
Iteration from \(i_0\) shows that
\(R_i\le C\theta_2^{i-i_0}+C/i^2\).  For \(i\ge J/2\) the
first term is at most \(C/i^2\), uniformly for all later \(i\),
when \(s\) is large.  Thus \(R_i=O(i^{-2})\) and
\(d_i=O(i^{-1})\).

Repeat the readout comparison with these improved bounds.  Each
nonlinear and remainder error is now \(O(i^{-2})\).  Consequently
the two unbalanced scalar coordinates agree with those of the
canonical history to \(O(i^{-2})\).  Each coordinate change to
\(W,Q\) differs from its linear part by \(O(i^{-2})\), by
Proposition~\ref{prop:map}.  The canonical critical estimates
therefore give \eqref{eq:traj-late-critical}.
\end{proof}

\subsection{A temperature secant at the matching index}

Let
\[
 \eps=\bc-b>0,\qquad \delta=\sqrt\eps,
 \qquad s\asymp\eps^{-1/3},\qquad
 J=\lfloor C_*\log s\rfloor,
\]
where \(s\) is chosen admissibly.  Consecutive admissible sides have
a fixed ratio, so the two-sided bound implicit in \(\asymp\)
is independent of \(\eps\).  These choices give
\begin{equation}\label{eq:traj-scale-slack}
 \eps s^2\longrightarrow0,
 \qquad \delta^{-1}=o(s^2),
 \qquad J=o(\delta^{-1}).
\end{equation}
The first inequality places \(b\) inside the analytic entry disk;
the second keeps every stop of order \(\delta^{-1}\) within the
cutoff range.

\begin{lemma}[Matching the temperature secant]\label{lem:traj-secant}
For the history with observation side \(s\),
\begin{equation}\label{eq:traj-secant-match}
 \frac{W_J(\bc-\eps)-W_J(\bc)}{-\eps}
      =J\bigl(\mathcal Pq_u+o(1)\bigr),
 \qquad
 \left\|\frac{Q_J(\bc-\eps)-Q_J(\bc)}{\eps}\right\|\le C.
\end{equation}
The same uniform bound holds for the secant of each unbalanced torus
remainder at stage \(J\) behind a stop of order \(\delta^{-1}\).
\end{lemma}

\begin{proof}
Entry analyticity and propagation in the small tube first give,
through stage \(J\), the crude bound
\begin{equation}\label{eq:traj-crude-proximity}
 \text{distance between the states at \(\bc-\eps\) and \(\bc\)}
       \le C\eps s^2e^{2\gamma J}.
\end{equation}
This statement uses each history's own norm.  The torus bound follows
by contracting its remainder with the same scalar inputs.  Since
\(\eps\asymp s^{-3}\) and \(2\gamma C_*<1\), the right side
decays as a negative power of \(s\); in particular it is smaller
than every fixed negative power of \(J\).  Induction with a fixed
margin to the boundary of the tube justifies the propagation up to
\(J\).

Denote the scalar secants by \(\widetilde V_i\), with denominator
\(-\eps\).  On \([J/2,J]\), Lemma~\ref{lem:traj-late-critical}
and \eqref{eq:traj-crude-proximity} allow us to integrate the
analytic derivatives of the map along the segment between its two
input states.  This gives the same form of scalar and remainder
inequalities as \eqref{eq:traj-linear-response} and
\eqref{eq:traj-derivative-contraction}.  The extra coefficients
from \eqref{eq:traj-crude-proximity} are exponentially small in
\(J\), and hence fit inside their displayed errors.  The crude
bound for any initial remainder secant is \(Cs^2e^{2\gamma J}\).
The weighted estimates \eqref{eq:traj-weighted-forward}--
\eqref{eq:traj-weighted-backward}, now on a fixed-fraction interval
ending at \(J\), therefore give
\begin{equation}\label{eq:traj-secant-slaving}
 \widetilde r_J\le\frac CJ|\widetilde V_J|+e^{-cJ}.
\end{equation}
Indeed the homogeneous contribution is bounded by a constant times
\(r^{c_1J}s^2e^{2\gamma J}\).  The choices of \(C_*\) and
\(\gamma\) make its exponent strictly negative; taking, for
example, an interval of length \(J/4\) within \([J/2,J]\)
suffices.  The estimate is uniform for the torus remainders behind
later stops.

To identify this secant, choose a canonical history with
\(s_k\asymp J\) and \(J\) in its fixed interior range.  In that
history we can compare the secant directly with its derivative.
For all stages through \(J\), analytic continuation in the
parameter is valid and bounded on a disk of radius
\[
 \rho_J=cJ^{-2}e^{-2\gamma J}
\]
about \(\bc\).  To see this, the entry displacement is bounded
by \(Cs_k^2\rho_J\), and its amplification through \(J\)
is bounded by \(Ce^{\gamma J}\); the resulting displacement
tends to zero.  The same argument and contraction control the
separate torus states.  Cauchy's estimate bounds second parameter
derivatives of the stopped states by \(C\rho_J^{-2}\).
Since \(\log\eps=-3J/C_*+O(1)\),
\begin{equation}\label{eq:traj-cauchy-error}
 \eps\rho_J^{-2}
 \le CJ^4\exp\{(-3/C_*+4\gamma)J\}
 \longrightarrow0.
\end{equation}
Thus the canonical scalar secant is
\(V_J^{(k)}+o(1)\), with analogous absolute control of its
remainder secants.  This is the only use of a second temperature
derivative; it is taken in the history with entry side comparable
to \(J\).

Compare the two histories through their physical readout secants.
Lemma~\ref{lem:grid-comparison} and
Lemma~\ref{lem:cutoff-derivative} apply uniformly on the real
parameter segment \([\bc-\eps,\bc]\), so their first-derivative
bounds apply also to the divided differences.  Taylor expansion of
\(\mathscr O_{J,s}\) from \eqref{eq:traj-normalized-diagnostic}
on the segment between each history's two temperature states,
followed by \eqref{eq:traj-secant-slaving}, gives
\begin{equation}\label{eq:traj-secant-comparison}
 D_J\widetilde V_J
   =D_JV_J^{(k)}
       +O\!\left(
           \frac{|\widetilde V_J|+|V_J^{(k)}|}{J}\right)+o(1),
\end{equation}
with the same uniformly invertible diagonal matrix as in
\eqref{eq:traj-response-readout}.  The multipliers and observation
errors are exponentially small because the first condition in
\eqref{eq:traj-parameter-choice} also gives
\(n_J^2(s/n_J)^4\le e^{-cJ}\).

Proposition~\ref{prop:response} gives \(|V_J^{(k)}|=O(J)\).
Absorbing \(C|\widetilde V_J|/J\) in
\eqref{eq:traj-secant-comparison} shows
\(|\widetilde V_J|=O(J)\), then
\(\widetilde V_J-V_J^{(k)}=O(1)+o(1)\).
Proposition~\ref{prop:response} and
\eqref{eq:traj-secant-slaving} prove all the assertions.
\end{proof}

\subsection{The nonlinear passage time}

We now follow the real perturbed trajectory beyond the matching
index.  The result includes exactly the two physical variance facts
needed in the following section.  Recall the microscopic torus
cosine \(H_n\) from \eqref{eq:physical-cosine}, and
\(\mathfrak g=1/(8\pi^2)\).

\begin{proposition}[Passage time and the variance diagnostic]
\label{prop:passage}
Put
\[
 k_0=\sqrt{3\mathcal P},\qquad T=\frac\pi{k_0}.
\]
For \(b=\bc-\eps\), use an admissible observation side
\(s\asymp\eps^{-1/3}\).  On every compact subinterval of
\((0,T)\), the scalar trajectory satisfies, uniformly for
\(w=\delta j\) in that subinterval,
\begin{equation}\label{eq:traj-scaled-limit}
 \frac{x_j(b)}\delta\longrightarrow k_0\cot(k_0w),
 \qquad
 \frac{y_j(b)}\delta\longrightarrow k_0\csc(k_0w),
 \qquad \|Q_j(b)\|=O(\delta^2).
\end{equation}
In particular, for every sequence \(\eps\downarrow0\) and
\(\delta j\to w\in(0,T)\),
\begin{equation}\label{eq:traj-before-variance}
                \Var_{b,n_j}^{\mathrm{tor}}H_{n_j}
                       \longrightarrow a\mathfrak g.
\end{equation}
For every sufficiently small fixed \(d>0\), there are exit indices
\(j_e=j_e(\eps,d)\) such that
\begin{equation}\label{eq:traj-exit}
 \delta j_e\longrightarrow T,\qquad
 x_{j_e}=-d+O(d^2),\qquad y_{j_e}=O(d),\qquad
 \|Q_{j_e}\|+\|\mathcal R_{j_e|j_e}^{\mathrm{tor}}\|=O(d^2).
\end{equation}
All states through the exit lie in the small-activity map domain.
For some \(c_d>0\),
\begin{equation}\label{eq:traj-exit-variance}
 \limsup_{\eps\downarrow0}
       \Var_{b,n_{j_e}}^{\mathrm{tor}}H_{n_{j_e}}
                 \le a\mathfrak g-c_d.
\end{equation}
The torus law in these assertions is the original periodic Bessel
height law in the zero-winding sector, with its common height
translation factored out.
\end{proposition}

\begin{proof}
Set
\(\Delta W_i=W_i(\bc-\eps)-W_i(\bc)\) and define
\(\Delta Q_i\) similarly.  We first prove that, for a sufficiently
small fixed \(\eta>0\), uniformly on
\(J\le i\le\eta/\delta\),
\begin{equation}\label{eq:traj-early-perturbation}
 \Delta W_i=-\eps i\bigl(\mathcal Pq_u+O(\eta^2)+o(1)\bigr),
 \qquad \|\Delta Q_i\|\le C\eps.
\end{equation}
Here the constant in \(O(\eta^2)\) is independent of sufficiently
small \(\eta\), and \(o(1)\to0\) as \(\eps\to0\).
The critical history has the estimates of
Lemma~\ref{lem:traj-late-critical} through every fixed multiple of
\(\delta^{-1}\), by \eqref{eq:traj-scale-slack}.

Bootstrap the estimate \(|\Delta W_i|\le C_1\eps i\).
The differentiated remainder bounds in
Proposition~\ref{prop:map} give
\[
 \|\Delta Q_{i+1}\|
 \le\theta_1\|\Delta Q_i\|
       +C_{C_1}\bigl(i^{-1}\eps i+(\eps i)^2\bigr).
\]
The second term is \(O_{C_1}(\eps)\) on the indicated range.
Lemma~\ref{lem:traj-secant} and contraction therefore give
\(\|\Delta Q_i\|\le C_{C_1}\eps\).
Subtract the two scalar recurrences and linearize their leading
quadratic terms about the critical trajectory.  The result is
\begin{equation}\label{eq:traj-difference-recurrence}
 \Delta W_{i+1}
       =\left(I+\frac{\mathsf M}{i}\right)\Delta W_i+e_i,
 \qquad
 |e_i|\le C_{C_1}\left(
          \frac{\eps\log i}{i}+\frac\eps i+\eps^2i^2\right).
\end{equation}
The first error uses the \(O(\log i/i^2)\) critical correction
and the differentiated cubic terms.  The second uses
\(\|\Delta Q_i\|=O(\eps)\) and the \(O(i^{-1})\) scalar
remainder row.  The last is the quadratic term in \(\Delta W_i\).

In the eigenbasis of \(\mathsf M\), the unperturbed products from
\(h\) to \(i\) have norm at most \(Ci/h\), and send \(q_u\)
to \((i/h)q_u\) exactly.  Variation of constants in
\eqref{eq:traj-difference-recurrence} and
Lemma~\ref{lem:traj-secant} therefore gives, after division by
\(\eps i\), a total error bounded by
\[
 o(1)+C_{C_1}\sum_{h\ge J}\frac{\log h+1}{h^2}
           +C_{C_1}\eps\sum_{h\le i}h
 \le o(1)+C_{C_1}\left(\frac{\log J}{J}+\eta^2\right).
\]
Choose \(C_1\) larger than the leading bound, then choose
\(\eta\) small.  The improved inequality closes the bootstrap.
The trajectory stays within the chosen tube throughout this
argument, by the same bound and the critical estimates.
This proves \eqref{eq:traj-early-perturbation}.  The unbalanced
remainder on every torus behind a stop of order \(\delta^{-1}\)
satisfies its corresponding bounds by the same contraction with
copied scalars.

For the continuous-scale limit, linearly interpolate
\(x_j/\delta,y_j/\delta\) at times \(w=\delta j\).
On every compact subinterval of \((0,\eta]\),
\eqref{eq:traj-early-perturbation} and the critical estimates bound
these interpolants uniformly and give \(Q_j=O(\delta^2)\).
The map in Proposition~\ref{prop:map} then gives bounded slopes and,
uniformly on such a subinterval, the limiting equations
\begin{equation}\label{eq:traj-ode}
                        X'=-Y^2,\qquad Y'=-XY.
\end{equation}
Indeed the scalar map error is \(O(\delta^3)\) there, while a
step has length \(\delta\) and the scalar coordinates were divided
by \(\delta\).  Its contribution to the slope tends to zero.
Compactness and integration of the discrete equations therefore
show that every subsequential limit solves
\eqref{eq:traj-ode}.

Applying \eqref{eq:traj-early-perturbation} with arbitrarily small
\(\eta\) determines its behavior at the left endpoint:
\begin{equation}\label{eq:traj-left-asymptotic}
 (X(w),Y(w))
   =\frac{(1,1)}w
       -w\bigl(\mathcal Pq_u+O(w^2)\bigr)
                  \qquad(w\downarrow0).
\end{equation}
The equations preserve \(X^2-Y^2\), whose value is
\(-3\mathcal P\) by \eqref{eq:traj-left-asymptotic}.
Thus \(X'=-(X^2+3\mathcal P)\).  Its pole at zero and the
positive sign of \(Y\) near zero give uniquely
\[
 X(w)=k_0\cot(k_0w),\qquad Y(w)=k_0\csc(k_0w).
\]
Every subsequential limit is the same, so convergence holds on
compact subintervals of \((0,\eta]\).

We extend it to each compact subinterval of \((0,T)\).
Starting from any fixed time in \((0,\eta)\), choose a bound
strictly larger than the supremum of the displayed solution on
the target compact interval.  While the discrete scalars divided
by \(\delta\) remain within this bound, remainder contraction
gives \(\|Q_j\|=O(\delta^2)\), so the scalar error in the Euler
approximation tends uniformly to zero.  The quadratic vector field
is Lipschitz on this bounded set.  The discrete Gronwall inequality
then gives uniform convergence up to the end of the interval and
excludes a first exit from the larger bound.  This proves
\eqref{eq:traj-scaled-limit}, including the analogous torus
remainder bounds behind these stops.

It remains to control the final part of the trajectory, where the
scaled solution diverges but the unscaled activity is still small.
Fix a time \(w_0\in(T/2,T)\).  Starting at
\(j_{\mathrm{start}}=\lfloor w_0/\delta\rfloor\), let \(j_e\) be the first
index with \(-x_{j_e}\ge d\), provided such an index exists.
We prove its existence and the estimates in
\eqref{eq:traj-exit} using a cone on which both scalar coordinates
have a definite direction of motion.

Take a large fixed \(M>k_0\) and let \(w_M\in(T/2,T)\) be the
outgoing time at which \(Y(w_M)=M\).  At an index with
\(\delta j\to w_M\), compact convergence gives
\[
 y_j=(M+o(1))\delta,
 \qquad -x_j/y_j=1+O(M^{-2})+o(1).
\]
The remainder estimate at this index can be made
\begin{equation}\label{eq:traj-cone-initial}
 \|Q_j\|+\|\mathcal R_{j|h}^{\mathrm{tor}}\|
                     \le K M^2\delta^2
\end{equation}
with \(K\) independent of large fixed \(M\).  Here \(h\) can
be any later stop up to a fixed sufficiently large multiple of
\(\delta^{-1}\).  To verify the independence, use the fixed
scaled interval on which the limiting outgoing \(Y\) increases
from \(M/2\) to \(M\).  Throughout it \(|W|\le CM\delta\)
for all sufficiently small \(\delta\), with an absolute \(C\),
and its number of discrete steps tends to infinity.  Contraction
bounds the forced remainder by \(CM^2\delta^2\) and removes
the contribution from the earlier state.  The threshold for
\(\delta\) may depend on \(M\), as is appropriate when \(M\)
is fixed before taking \(\delta\to0\).

For sufficiently small coordinates consider
\begin{equation}\label{eq:traj-outgoing-cone}
 y>0,\qquad \frac12\le q:=-x/y\le2,
 \qquad \|Q\|+\|\mathcal R^{\mathrm{tor}}\|\le K_1y^2.
\end{equation}
Choose \(K_1\) sufficiently large for the remainder contraction.
In this cone the scalar map gives
\begin{align}
 y_+&=y\{1+qy+O_{K_1}(y^2+j^{-2}y)\},
       \label{eq:traj-y-growth}\\
 q_+-q&=y(1-q^2)+O_{K_1}(y^2+j^{-2}y).
       \label{eq:traj-ratio-growth}
\end{align}
At each endpoint of \([1/2,2]\) the leading term in
\eqref{eq:traj-ratio-growth} points inward.  It also points
uniformly inward on fixed small neighborhoods of the endpoints.
Away from those neighborhoods a sufficiently small step cannot
cross either endpoint.  Thus the ratio condition is preserved
for small \(y\) and large \(j\).  Equation~\eqref{eq:traj-y-growth}
gives constants \(0<c<C\) such that
\[
                    y+cy^2\le y_+\le y+Cy^2.
\]
Remainder contraction preserves the last condition in
\eqref{eq:traj-outgoing-cone}: its forcing is at most \(C'y^2\),
so choose \(K_1>C'/(1-\theta_1)\), and use \(y_+\ge y\).
This applies to the plane remainder and separately to every
stopped torus under consideration.

For large fixed \(M\), \eqref{eq:traj-cone-initial} places the
trajectory in this cone.  Choose \(d\) small enough that the cone
estimates hold until the first crossing of \(-x=d\), including
that crossing.  Before it, \(y\le2d\), and reciprocal summation
gives
\[
 \frac1{y_i}-\frac1{y_{i+1}}
   =\frac{y_{i+1}-y_i}{y_iy_{i+1}}\ge c'>0.
\]
It follows that the exit occurs within \(C/(M\delta)\) further
steps: otherwise the last inequality would force \(1/y_i<0\).
At the crossing, the change of \(x\) is \(O(y^2)=O(d^2)\).
The cone and remainder bounds therefore give the last three
assertions in \eqref{eq:traj-exit}.

Compact convergence on \([w_0,T-\zeta]\), for any fixed
\(\zeta>0\), shows that no such exit occurs there when
\(\delta\) is small.  Hence \(\liminf\delta j_e\ge T\).
The preceding upper bound gives
\[
 \limsup_{\delta\downarrow0}\delta j_e\le w_M+C/M.
\]
Since \(w_M\uparrow T\) as \(M\to\infty\), this proves
\(\delta j_e\to T\).  Throughout this argument all unscaled
coordinates remain small; no activity map is used beyond its
domain.  The fixed horizon of stopped tori can be taken larger
than \((T+1)/\delta\), and the exit lies before it for small
\(\delta\).

Finally apply the variance readout of
Proposition~\ref{prop:readouts}.  At a stop with
\(\delta j\to w\in(0,T)\), the scalar activity is \(O(\delta)\)
and the remainder is \(O(\delta^2)\), so the variance converges
to \(a\mathfrak g\).  At the exit, the coordinate change gives
\[
 t_{j_e}=-\frac d{b_*}+O(d^2),\qquad
 \|\mathcal R_{j_e|j_e}^{\mathrm{tor}}\|+\nu_{j_e}^2=O(d^2).
\]
Here \(\nu_{j_e}\) is the full stopped activity size.  The
readout multiplier tends to \(a\mathfrak g\), and therefore
\[
 \Var_{b,n_{j_e}}^{\mathrm{tor}}H_{n_{j_e}}
     =a\mathfrak g-\frac{a\mathfrak g}{b_*}d
                        +O(d^2)+o_{\eps\downarrow0}(1).
\]
To pass from the cutoff observation to \(H_n\), remove the
independent noise and use the sampling comparison and cutoff
bound.  Their errors tend to zero both before and at the exit,
since \(\log n_j=O(\delta^{-1})=o(s^2)\) and \(s/n_j\to0\).
Taking \(d\) sufficiently small gives
\eqref{eq:traj-exit-variance}, for example with
\(c_d=a\mathfrak g d/(2b_*)\) after decreasing its allowed range.
\end{proof}
